\documentclass[11pt]{report}
\usepackage[margin=1in]{geometry}
\usepackage[utf8]{inputenc}
\usepackage{amsmath,amssymb}
\usepackage{graphicx}
\usepackage{booktabs}
\usepackage{caption}
\usepackage[hidelinks]{hyperref}
\usepackage{enumitem}

\newcommand{\dkl}{$\Delta$KL}

\title{\bfseries Trademark resonance and business-model diffusion\\
\large An integrated report \\[0.3em]
\normalsize Part I: Construct validity of $\Delta$KL \quad\textbar\quad Part II: Phase-0/1 of the diffusion program}
\author{Eric Silver}
\date{June 2026 --- working report}

\begin{document}
\maketitle
\tableofcontents
\clearpage

% ========================================================================
\chapter{What does \dkl{} measure? Construct validity of the trademark resonance signal}
\noindent\textit{Within-firm patent co-movement and external-list discriminant validity.}\medskip

\noindent\textbf{Abstract.}\quad For each USPTO trademark filing we score
$\Delta\mathrm{KL}=\mathrm{KL}(\text{filing}\,\|\,\text{past})-\mathrm{KL}(\text{filing}\,\|\,\text{future})$
on its goods/services text against same-class reference windows --- DeDeo
\emph{resonance}: positive when a filing's vocabulary is novel against the recent
past yet close to the near future (``the field moved toward it''). Does this
measure \emph{innovation}? Using the public-firm panels, three results sharpen
the construct. (1)~\dkl{} is essentially uncorrelated with patents in the cross
section (Pearson $+0.010$, Spearman $-0.015$) and, within a firm over time, a
\emph{mild substitute} for patenting ($-0.05\sigma$ per log-patent, $t\!\approx\!-8$,
strongest at the one-year patent lag) --- robust to filing-count and patent-matching
choices. (2)~Yet firms that BCG (2021--23) and MIT Technology Review (2015)
independently rank ``most innovative'' score $+0.29$ to $+0.50$ population-$\sigma$
\emph{higher} on \dkl{} ($p<0.012$; pooled $+0.34\sigma$, $p=0.0003$). (3)~The
headline ``$+28$pp per $\sigma$'' four-year debut-return effect is a small-sample
right-tail artifact: its horizon profile is non-monotone, peaking at 4y and
collapsing at 5y, and it runs $2.4\times\!\to\!9.3\times$ the firm-year estimate as
the horizon lengthens. Together these say \dkl{} is \emph{not} a proxy for
patent-style technical invention but \emph{does} track market/business-model
innovation as recognized by expert observers. The right question is not
``innovation versus churn'' but \emph{which} innovation construct \dkl{} names.

\section{The construct question}
The program's theory (PROPOSAL.md) defines innovation as \emph{novelty that
diffuses}: a business concept is ``innovative'' if its elements are
\emph{replicated} after the firm introduces them. The resonance statistic
operationalizes exactly this --- high \dkl{} marks language that was unusual at
filing and then adopted. On that theory, a filing being ahead of a class whose
vocabulary subsequently moves toward it is not a churn \emph{confound} to be
purged; it \emph{is} the innovation signal. The sharper question is therefore not
whether to subtract ``category churn'' out of \dkl{}, but whether \dkl{} names a
recognizable innovation \emph{construct} and how it relates to the established
ones. We test it against the two best external yardsticks available in the
repository: patents (technical invention) and expert ``most-innovative'' lists
(market/business-model innovation).

\paragraph{Scope.} The raw 10M-filing corpus is not on this branch, so the
corpus-level tests (term provenance, half-life sensitivity, ID-Manual
codification, attorney fixed effects) are deferred (\S\ref{sec:blocked}). All
results below run on the published firm-year panels and the external-list
crosswalk.

\section{Within-firm: \dkl{} and patenting are temporal substitutes}
Panel: 174{,}569 firm-years for 14{,}463 trademark owners name-matched to a
PatentsView assignee, collapsed to one row per (firm, year). We regress
standardized \dkl{} on $\log(1+\text{patents})$ with firm and year fixed effects
and firm-clustered standard errors; $\sigma(\Delta\mathrm{KL})=0.214$.

\begin{table}[h]\centering\small
\caption{Within-firm \dkl{} vs.\ patenting. Coefficients in \dkl{} standard
deviations per unit $\log(1+\text{patents})$.}
\begin{tabular}{lrrr}
\toprule
Specification & coef ($\sigma$) & SE & $t$ \\
\midrule
Between-firm correlation (firm means) & $-0.020$ & --- & --- \\
Within-firm correlation (firm+year demeaned) & $-0.015$ & --- & --- \\
Two-way FE, contemporaneous & $-0.048$ & $0.006$ & $-7.9$ \\
\quad robustness: patents = max over matches & $-0.048$ & $0.006$ & $-8.0$ \\
\quad controlling $\log(\text{n\_filings})$ & $-0.051$ & $0.006$ & $-8.4$ \\
Lead/lag: patents$_{t-1}$ & $-0.054$ & $0.009$ & $-5.8$ \\
\quad\quad\quad\ \ patents$_{t}$ & $-0.019$ & $0.010$ & $-1.9$ \\
\quad\quad\quad\ \ patents$_{t+1}$ & $+0.010$ & $0.010$ & $+1.0$ \\
\bottomrule
\end{tabular}
\end{table}

The coefficient is negative and precisely estimated, and \emph{strengthens} when
we control for the number of filings (which enters positively, $+0.055$,
$t=10.5$) --- ruling out the mechanical worry that high-patent years simply have
more filings whose averaging shrinks the mean toward the class. The effect is
largest at the one-year patent \emph{lag}: a burst of patenting is followed by
\emph{lower} trademark resonance the next year. Measurement noise from the fuzzy
owner$\to$assignee join attenuates toward zero, so the true within-firm relation
is, if anything, more negative.

\begin{figure}[h]\centering
\includegraphics[width=0.62\textwidth]{results/wsC_leadlag.png}
\caption{Within-firm \dkl{} declines with recent patenting; no positive
co-movement at any lead or lag.}
\end{figure}

\paragraph{Reading.} This decisively refutes \emph{convergent} validity with
patents --- not merely the near-zero cross-sectional correlation, but a genuine
within-firm \emph{opposition}. In the years a firm leans on patents it writes
more class-typical trademark language; in the years it pioneers commercial
vocabulary it patents less. Patents and \dkl{} capture different, partly
substituting modes of novelty.

\section{Discriminant validity against expert innovation lists}
If \dkl{} names an innovation construct at all, firms an independent panel calls
``most innovative'' should sit above the population. We compare the firm-mean
\dkl{} of list members (matched to a CIK in our public panel) against the
non-listed baseline (population: 7{,}431 firms, mean $-0.018$, $\sigma=0.175$).

\begin{table}[h]\centering\small
\caption{\dkl{} of expert-rated innovators vs.\ the non-listed baseline.}
\begin{tabular}{lrrrr}
\toprule
List & $n$ matched & listed mean & diff (pop-$\sigma$) & Mann--Whitney $p$ \\
\midrule
BCG Most Innovative (2021--23) & 31 & $+0.032$ & $+0.285$ & $0.012$ \\
MIT TR50 Smartest (2015)       & 17 & $+0.069$ & $+0.497$ & $0.0006$ \\
Pooled (union)                 & 43 & $+0.042$ & $+0.344$ & $0.0003$ \\
\midrule
\multicolumn{5}{l}{\textit{Convergent reference --- patents (174{,}569 firm-years):}}\\
Pearson (raw / log) & & $+0.010 / -0.009$ & & \\
Spearman (rank)     & & $-0.015$ & & $8\!\times\!10^{-10}$ \\
\bottomrule
\end{tabular}
\end{table}

\begin{figure}[h]\centering
\includegraphics[width=0.62\textwidth]{results/wsD_discriminant.png}
\caption{Population distribution of firm-mean \dkl{} with the BCG and MIT TR50
list means overlaid. Expert-rated innovators sit near the 67th--76th percentile.}
\end{figure}

\paragraph{Reading, with caveats.} \dkl{} aligns with expert, non-patent,
non-financial innovation judgments while being orthogonal-to-negative with
patents. That is the signature of a \emph{market/business-model} innovation
construct rather than a technical-invention one. Caveats are real and limit this
to suggestive: $n$ is small (only the US-public, CIK-matchable slice of global
lists matched), the comparison is each firm's all-years mean against a recent
list (no temporal alignment), and survivorship inflates the listed group. It is
nonetheless the first \emph{out-of-corpus} validation in the program --- the
lists were built without trademark text.

\section{The +28pp return effect is an artifact --- drop it}
The predictor is $z$-scored in \texttt{returns\_regression.py}, so the
$+0.284$ debut coefficient really is $+28$pp per $\sigma$. But it is not credible
as a headline: the debut excess-return profile is non-monotone and the debut /
firm-year ratio grows monotonically with horizon --- the fingerprint of a few
right-tail winners in a 1{,}447-firm sample, not a return premium.

\begin{table}[h]\centering\small
\caption{Excess-return coefficient on \dkl{} (per $\sigma$), by horizon.}
\begin{tabular}{lrrr}
\toprule
Horizon & Debut ($n\!\approx\!1.4$k) & Firm-year ($n\!\approx\!16$k) & ratio \\
\midrule
1y & $+0.036\ (t\,2.0)$ & $+0.015\ (t\,3.0)$ & $2.4\times$ \\
2y & $+0.067\ (t\,2.5)$ & $+0.035\ (t\,3.8)$ & $1.9\times$ \\
3y & $+0.192\ (t\,3.5)$ & $+0.050\ (t\,3.7)$ & $3.8\times$ \\
4y & $+0.284\ (t\,3.3)$ & $+0.041\ (t\,2.5)$ & $6.9\times$ \\
5y & $+0.149\ (t\,2.3)$ & $+0.016\ (t\,0.8)$ & $9.3\times$ \\
\bottomrule
\end{tabular}
\end{table}

\noindent Report the firm-year association ($+1.5$ to $+5$pp per $\sigma$, decaying
to insignificance by 5y) as the defensible financial fact, pending value-weighted,
delisting-adjusted, $5/95$-winsorized re-estimation of the debut panel.

\section{Discussion: which innovation construct?}
The ``innovation vs.\ churn'' dichotomy dissolves. \dkl{} is not a noisy patent
proxy contaminated by category churn; it is a measure of a \emph{different}
construct that happens to be near-orthogonal to patents. The evidence:

\begin{itemize}[nosep]
\item \dkl{} $\perp$ patents (cross-section) and $<0$ within-firm $\Rightarrow$
it does not track technical invention;
\item \dkl{} $> 0$ for BCG/MIT innovators $\Rightarrow$ it tracks market/
business-model innovation as expert observers perceive it.
\end{itemize}

\noindent On the ladder of names: the strong claim ``\dkl{} is \emph{the}
innovation measure'' is not licensed. The defensible primitive is \emph{lexical
resonance} --- a filing's forward-leaning position in its class vocabulary --- and
its validated economic content is \emph{market-facing innovation as a construct
distinct from invention}. The diffusion-as-innovation theory survives; what fails
is the assumption that this construct should converge with patents. It should
not, and does not.

\paragraph{Bridge to the diffusion program (Part II).} The Phase-0 test of
\texttt{PROPOSAL\_diffusion.md} on the four-class ``Uber-for-X'' subset
\{009, 042, 039, 035\} (2.32M granted filings, Part~II of this report) passes
both §5 kill conditions: the date/class shuffle null is rejected ($z=+46$ on
the lag-1 autocorrelation of LDA-theme trajectories), independently replicated
on a separately-fitted NMF panel ($z=+60.7$), and signed-ahead-precedes-behind
is directionally correct ($61\%$ of $(\theta, c)$ cells negative, modest
magnitude consistent with the $0.97$ collinearity between prospective and
retrospective KL noted above). The diffusion register --- themes as the unit,
not the per-filing scalar --- is the live continuation of the construct argument
made here.

\paragraph{The localized innovation signal materializes (RQ6 in Part II).}
The diffusion program's RQ6 --- does borrowed novelty pay at the firm level? ---
is the test this chapter said would be the right place to look for an
innovation reading. With the SEC Financial Statement Data Sets crosswalked to
USPTO owners (7{,}233 exact matches, 4{,}999 CIK-linked firms in the
borrowed-novelty panel), the firm-year two-way FE regression of financial
outcomes on within-firm borrowed-gap finds: \textbf{contemporaneous gross margin
$+1.4$pp per $\sigma$ ($t=2.1$)} that fades by $t+1$, and \textbf{firms with
higher borrowed novelty have markedly lower R\&D intensity} cross-sectionally.
That mirror image --- borrowed novelty buys some margin in the year of the
filing while being a within-firm substitute for technical R\&D --- coheres
exactly with the within-firm $\Delta$KL/patent substitution result in §1.2 above:
borrowed novelty is the market-facing innovation channel; the firms who lean
on it lean less on the technical-invention channel. The localized innovation
reading is licensed on this margin-and-substitution shape, not on a global
convergent claim.

\section{Within-class diagnostics on a single technology corpus (Class 009)}

The firm-year results above use the aggregated panel. To stress-test the
construct's robustness at the filing level, we ran the full DeDeo flat-window
pipeline on a complete Class 009 (Software \& Electronics) corpus: 1{,}807{,}636
filings, 1893--2026, vocabulary $V = 151{,}193$ (min $df = 50$), per-filing
prospective + retrospective KL at $W \in \{3, 5, 7\}$ years. Four diagnostics
matter for what we can claim.

\paragraph{Pros/retr collinearity.} Across 1.56M scored filings,
$\mathrm{corr}(\mathrm{pros\_kl}, \mathrm{retr\_kl}) = \mathbf{+0.971}$. Signed
$\Delta$KL is a \emph{small residual of two large, highly correlated quantities}.
That is the point of the difference (it isolates timing from raw
distinctiveness) but also the fragility (small residuals are noisy at the
filing level).

\paragraph{Length artifact (F0b).} Mean $|\Delta\mathrm{KL}|$ falls
monotonically with filing length: $0.293$ for $<5$ terms, $0.265$ for 5--14,
$0.205$ for 15--39, $0.147$ for $\geq 40$. Short filings carry inflated
magnitudes by construction. Any per-filing analysis should down-weight or
control for length.

\paragraph{Window sensitivity (F0c).} $\Delta$KL \emph{magnitudes} are kernel-
sensitive, but the \emph{sign} is robust: only $10.4\%$ of filings flip sign
between $W=5$ and $W=3$ years and $6.9\%$ between $W=5$ and $W=7$, well under
the 20--30\% fragility threshold. The resonance direction (ahead vs.\ behind)
is not a kernel artifact.

\paragraph{Burn-in.} Mean signed $\Delta$KL at 1990 is $+2.05$ because the
pre-1990 software past is thin; from 1995 onward the per-year mean settles
within $\pm 0.07$ of zero. Restrict pre-1995 to a burn-in window.

\paragraph{Face validity post-burn-in.} Among 1996--2020 filings with $\geq 10$
terms, the top signed-$\Delta$KL marks are dominated by \textbf{2018
blockchain-era language}: STRAINCHAIN, CORPCOIN, MEDEXCHAIN, THE BLOCKCHAIN
BANKER, MEDXCHAIN, ART DEVICE. These are filings whose vocabulary was novel
at filing time and adopted afterward --- the construct firing on exactly what
it claims.

\paragraph{Term provenance (within-class Workstream D).} For each of 872{,}747
post-burn-in filings (1996--2020) we computed \texttt{novel\_share} = fraction
of its tokens whose first appearance in Class 009 was within $K=3$ years prior.
$\mathrm{corr}(\texttt{novel\_share}, \Delta\mathrm{KL}) = \mathbf{+0.29}$
(Pearson) and $+0.21$ (Spearman, $p \approx 0$). Filings with $\sim 0$ recently-
introduced class tokens average $\Delta\mathrm{KL} = -0.03$; filings with $\sim
5\%$ new tokens average $+0.10$. \textbf{The resonance signal is genuinely
new-vocabulary-driven within the class}, not recombination of established
terms (the corpus-new vs class-new split needs cross-class data; here we can
only confirm new-to-class drives the signal). The cross-class version on the
full diffusion subset is covered in Part~II.

\section{What remains, and the decision rules}\label{sec:blocked}
These need the raw corpus (\texttt{data/processed/}) or external data not on this
branch. Each is scripted or specified; the decision rule is stated so the call is
made ex ante.
\begin{description}[nosep,style=nextline]
\item[Interaction estimand ($\beta\times$churn).] Needs per-NICE-class \dkl{}
and an \emph{independently constructed} churn measure (vocabulary turnover, not
the resonance kernel). Tests whether \dkl{}'s innovation content is moderated by
class dynamism. Rule: if the slope is positive only where churn $>0$, the
construct is intrinsically a firm$\times$class-phase quantity (consistent with the
diffusion theory), not a censored firm attribute.
\item[Term provenance (three-way).] Split \dkl{} into corpus-new, class-new
(recombination), and class-extant tokens. Rule: load on recombination $\Rightarrow$
rename the firm component ``recombinant novelty,'' not foresight.
\item[Half-life sensitivity.] Re-score sign of \dkl{} at $H\in\{1,2,4\}$y. Rule:
$>$20--30\% sign flips $\Rightarrow$ report \dkl{} as kernel-relative.
\item[ID-Manual codification (external).] Do high-\dkl{} filings use terms the
USPTO later standardizes, while those terms are still rare? This is the strongest
available external validation of the diffusion claim; needs the ID Manual version
history.
\item[Attorney fixed effects.] Needs an XML re-parse for correspondent of record.
Kills the ``drafting-fashion alone'' null; not a clean innovation extractor
(counsel sorting absorbs signal).
\item[External survival outcome.] Bankruptcy (PACER) or web-presence persistence
on a validation subsample, to break the trademark-maintenance circularity in the
survival U-shape.
\item[BISG demographics.] The ethnic-cluster work should replace surname-only
imputation with BISG using the owner addresses (verify the address is the owner's,
not the correspondent's, first).
\end{description}

% ========================================================================
\chapter{Business-model diffusion: Phase-0 and first-pass Phase 1}
\noindent\textit{Cross-class theme transit, LDA themes, §5 kill conditions, and a first read of RQ1--RQ4 / NMF / RQ5 / RQ6 in four NICE classes (009, 042, 039, 035).}\medskip

\noindent\textbf{Abstract.}\quad A first, descriptive test of
\texttt{PROPOSAL\_diffusion.md} on a four-class ``Uber-for-X'' subset ---
software (009), tech services (042), transport (039), advertising/retail (035),
2.35M granted filings 1990--2024. Three measurements. (1)~\textbf{Phrase-level
transit}: business-model vocabularies originate overwhelmingly in software/tech
and arrive in transport and retail years later (cloud +1--4y, AI +2y,
as-a-service +3--4y, mobile-app +3y, streaming +5--6y, AR +7y, blockchain +4y).
(2)~\textbf{The filing-level within/cross scalar gap} does \emph{not} detect
transit on its own --- high-within/low-cross filings are generic vocabulary in
an off-class, not coherent imports. (3)~A v0 LDA theme model ($T=50$, fit on a
200k stratified sample, transformed across the full 2.32M-filing pool)
\textbf{passes both} PROPOSAL §5 kill conditions: the shuffle null is rejected
decisively (real mean lag-1 autocorrelation of theme$\times$class prevalence
trajectories $+0.687$ vs.\ shuffle $-0.027$, $z=+46$), and the
signed-ahead-precedes-behind check is directionally correct ($61\%$ of
$(\theta, c)$ cells negative; mean Spearman
$\rho(\Delta\mathrm{KL},\,\text{year}) = -0.053$), modest in magnitude.
Verdict: the diffusion construct survives the gating tests; the theme is the
right unit (not the per-filing scalar); proceed to Phase 1. A first-pass
Phase~1 read on the 50-theme model: 46/50 themes cross the 1\% floor in at
least one class, with 12 reaching all four (RQ1, RQ2);
\textbf{class-to-class transit is asymmetric and software-led} --- 33 of the
between-class flows originate in software/tech, only 4 return there. Bass
fits ($R^2 \geq 0.7$) on 118 $(\theta, c)$ cells give median $p = 0.017$,
$q = 0.114$, $q/p \approx 7$, mostly platform-like (RQ3). Provenance among
the 50 earliest carriers per theme is 73.7\% entrant-side in raw terms (RQ4)
but the year-matched population baseline is \emph{higher} at 76.4\%, so the
mean theme excess is $-2.7$pp; the raw Schumpeter Mark~I read is artifactual
--- themes inherit the corpus's high overall debut rate without adding to it.
A NMF cross-method check ($T=50$, same vocabulary) replicates the diffusion
structure decisively (NMF shuffle null $z=+60.7$) and confirms 14--28\% of
topics as strongly cross-method-stable (Jaccard $\geq 0.3$--$0.5$) --- the
§2.4 firewall keeps only this subset for headline. The SEC Financial Statement
Data Sets panel (2009--2024) now also enables: RQ5 finds lagged class
\emph{revenue growth} predicts theme uptake (coef $+9.2$, $p=0.048$, $n=57$),
lagged margin does not --- direction supports Schmookler/demand-pull. RQ6
finds firms with high filing-year borrowed novelty earn modestly higher
contemporaneous gross margin ($+1.4$pp per $\sigma$, $t=2.1$) but with
\emph{lower} R\&D intensity --- a market-facing innovation signal that
\emph{substitutes} for technical invention rather than complementing it.

\section{Setup}
Four classes chosen to exercise the canonical cross-industry transit: a software
goods class (009), a tech-services class (042), and two ``physical/commercial''
destination classes (transport 039, advertising/retail 035). We condition on
\textbf{grant} (registration date present; §2.1 of \texttt{PROPOSAL\_diffusion.md}):
912k/575k/101k/765k granted filings respectively, 2{,}352{,}913 pooled overall,
with the kill-condition analyses restricted to filing years 1990--2024
(2{,}317{,}899 filings).

\section{Result 1 --- phrase-level theme transit (positive)}
For a curated set of multi-token business-model phrases (matched by regex on raw
goods/services text, so stopword removal does not mangle ``on demand'' /
``as a service''), we compute per (class, year) prevalence among granted filings
and record the first year each class crosses a 1\% floor.

\begin{table}[h]\centering\small
\caption{First year a theme reaches 1\% prevalence, by class. ``--'' = never.}
\begin{tabular}{lrrrr}
\toprule
Theme & software & tech-svcs & transport & advert/retail \\
\midrule
as-a-service           & 2011 & \textbf{2009} & 2013 & 2012 \\
mobile-app             & 2010 & \textbf{2009} & 2012 & 2012 \\
cloud                  & 2012 & \textbf{2010} & 2011 & 2014 \\
streaming              & \textbf{2008} & 2008 & 2013 & 2014 \\
artificial-intelligence& 2017 & \textbf{2016} & 2018 & 2018 \\
augmented/virtual-reality & \textbf{2015} & 2016 & 2022 & 2022 \\
blockchain / crypto    & 2021 & \textbf{2018} & --   & 2022 \\
real-time              & \textbf{2001} & 2003 & 2009 & 2014 \\
on-demand              & 2017 & 2016 & \textbf{2014} & 2017 \\
\bottomrule
\end{tabular}
\end{table}

\noindent Of the themes that clear the floor in a destination class, almost all
originate in software/tech (009/042) and arrive in transport/retail with lags of
1--13 years. The pattern is the proposal's central phenomenon, measured directly.
\textbf{Honest caveats.} The 1\% floor is sensitive to class size (039 is
$\sim$8$\times$ smaller, so its floor trips on fewer filings);
\texttt{peer-to-peer} never clears 1\% as an exact phrase; and \texttt{on-demand}
originates in \emph{transport} (2014), a genuine counter-example.

\section{Result 2 --- the scalar within/cross gap does not detect transit (negative)}
For each filing we computed prospective novelty against its own class's past
(\textit{within}) and against the pooled-corpus past (\textit{cross}), in a shared
116{,}783-term universal vocabulary. Within-novelty is typically \emph{below}
cross-novelty (own class predicts a filing better than the corpus), so the
``high-within/low-cross = borrowed'' cell is small (3--5\% by class) and its top
members --- EYEGLASSES DIRECT (009), GAP KIDS (035), EURO-FRAMES (035) --- are
generic vocabulary that is merely locally uncommon, not coherent business-model
imports. The scalar gap conflates ``locally rare generic words'' with ``imported
theme,'' precisely the vocabulary-migration-$\neq$-idea-migration threat (§4.2).
It is a useful \emph{filter} (transit candidates should be high-within/low-cross)
but not a \emph{detector} on its own.

\section{Result 3 --- LDA themes pass the §5 kill conditions}

\paragraph{Theme model.} A 50-topic LDA was fit on a stratified sample of 200{,}000
granted filings (50k per class, capped at class size; 039 smaller after sampling
constraints), with the resulting model used to transform the full 2.32M-filing
pool to a per-filing topic distribution. The vocabulary (unigram+bigram,
$\min\,df=200$) had $V=69{,}361$ terms; LDA fit took 5 min, full-corpus
transform 3 min.

\paragraph{Face validity.} The top-word lists are recognizable business-model
components --- e.g., T0 \{health, care, wellness, healthcare\}; T1 \{autonomous,
loading, unloading, truck, material handling\}; T2 \{store, retail, wholesale\};
T10 \{electronic transmission, computer networks, data, communication\}; T30
\{online, website, non-downloadable, information\}; T49 \{vehicles, traffic,
bicycle, lights\}. We make no claim that 50 unsupervised topics constitute the
final theme set --- §2.4's cross-method-stability and human-validation firewall
is still owed before any individual theme is treated as headline. The kill-
condition tests below are the cheap go/no-go check on whether \emph{any} theme
extraction can support the program.

\paragraph{Kill condition 1 --- shuffle null: pass.}
For each (theme $\theta$, class $c$) prevalence trajectory over 1990--2024, we
compute the lag-1 temporal autocorrelation. Real themes should have
\emph{positive} autocorrelation (smooth ramp-up of adoption); random
(class, year) label permutations should average near zero.

\begin{table}[h]\centering\small
\caption{Lag-1 autocorrelation of $(\theta, c)$ prevalence trajectories.}
\begin{tabular}{lcc}
\toprule
& mean autocorr & cells \\
\midrule
Real panel & $+0.687$ & 200 \\
Shuffle 1 & $-0.041$ & 200 \\
Shuffle 2 & $-0.041$ & 200 \\
Shuffle 3 & $-0.003$ & 200 \\
Shuffle 4 & $-0.024$ & 200 \\
\midrule
Shuffle mean $\pm$ sd & $-0.027 \pm 0.016$ & --- \\
\multicolumn{2}{l}{$z = (\text{real} - \text{shuffle})/\text{sd}$} & $\mathbf{+46}$ \\
\bottomrule
\end{tabular}
\end{table}

\noindent The real panel's themes ramp up smoothly within their classes; the
shuffled panel is noise around zero. The proposal §5 ``date/class shuffle null''
is rejected at any plausible significance level. The construct is not an
artifact of arbitrary panel slicing.

\paragraph{Kill condition 2 --- signed-ahead-precedes-behind: directionally pass, modest.}
For each $(\theta, c)$ cell with $\geq 50$ filings carrying $\theta$ above
$\tau = 0.20$ topic-weight, we compute Spearman $\rho$ between filing-level
signed $\Delta\mathrm{KL}$ (resonance) and filing year. Adoption-curve theory
predicts $\rho < 0$: early carriers are signed-\emph{ahead} (resonant; novel
language whose class then adopted it), late carriers signed-\emph{behind} (tired;
language the class has played out).

\begin{table}[h]\centering\small
\begin{tabular}{lc}
\toprule
$(\theta, c)$ cells & 197 \\
Mean $\rho(\Delta\mathrm{KL},\,\text{year})$ & $-0.053$ \\
Share with $\rho<0$ & $\mathbf{61.4\%}$ \\
\bottomrule
\end{tabular}
\end{table}

\noindent A majority of cells move in the predicted direction --- the
adoption-curve interpretation is not rejected --- but the per-cell effect is
small. This is consistent with the within-class diagnostics on Class~009 from
Part~I, §1.6: prospective and retrospective KL are 0.97 correlated, so
per-filing signed $\Delta\mathrm{KL}$ is a small residual and a noisy individual
adoption-position marker. The aggregate sign comes through; treating signed
$\Delta\mathrm{KL}$ as a clean per-filing curve coordinate is asking more than
it can do.

\section{First-pass Phase 1: RQ1--RQ4 on the 50-theme model}

\paragraph{RQ1 --- theme-arrival rate.} 46 of the 50 LDA topics cross the 1\%
prevalence floor in at least one class at least once; 4 remain niche
throughout. With the topic count fixed at $T = 50$ the proper Heaps' test is
deferred to a variable-$K$ pass, but as a within-fixed-$T$ check: the log-log
slope of cumulative-arrived-themes against cumulative filings is $\approx
+0.04$ --- the arrival is essentially front-loaded, most of the 46 themes are
already evident relatively early in the 1990--2024 window.

\paragraph{RQ2 --- breadth and class flow.}
\begin{table}[h]\centering\small
\caption{Theme breadth across the four classes.}
\begin{tabular}{cc}
\toprule
Classes reached & Themes \\
\midrule
1 (niche, class-specific) & 13 \\
2 & 6 \\
3 & 15 \\
4 (general-purpose)       & 12 \\
\bottomrule
\end{tabular}
\end{table}

\noindent 33 themes reach $\geq 2$ classes. The class-to-class transit
ordering (which class a theme arrives in second, after its origin) is
\emph{asymmetric and software-led}:
\begin{itemize}[nosep]
\item software $\to$ tech-svcs: 14;\quad software $\to$ transport: 11;\quad
software $\to$ advert/retail: 8 (33 transits originating in software);
\item transport $\to$ tech-svcs: 7;\quad transport $\to$ advert/retail: 6;
\item tech-svcs $\to$ software: 4 (the only meaningful return flow).
\end{itemize}
\noindent This confirms the phrase-level transit asymmetry of §2.2 at the
theme level: software is a net exporter of business-model vocabulary to the
other three classes. With only four nodes a centrality estimate is premature;
a 45-class build would let us locate each class in the business-model space.

\paragraph{RQ3 --- Bass $p/q/m$ on theme adoption.} For each $(\theta, c)$ cell
where the theme reached at least the floor and was tracked for $\geq 6$ years
we fit Bass(1969) to the normalized cumulative prevalence and retain fits with
$R^2 \geq 0.7$.

\begin{table}[h]\centering\small
\caption{Bass fits across $(\theta, c)$ cells.}
\begin{tabular}{lr}
\toprule
Fits retained ($R^2 \geq 0.7$) & 118 \\
median $p$ (innovation coef.) & $0.017$ \\
median $q$ (imitation coef.)  & $0.114$ \\
median $q/p$                   & $\sim 7$ \\
``platform-like'' ($q/p<20$)   & 88 \\
``fad-like'' ($q/p\geq 50$)    & 9 \\
\bottomrule
\end{tabular}
\end{table}

\noindent The bulk of the diffusion is platform-like (broad, gradual adoption,
moderate word-of-mouth); a small tail of $\sim 8\%$ shows fad-shaped curves
(very high $q/p$, sharp imitation pickup). The median $q/p \approx 7$ is in
the range typically reported for consumer-product Bass fits, lending external
plausibility to the theme$\times$class adoption curves as real adoption events.

\paragraph{RQ4 --- entrant vs.\ incumbent provenance.} For each theme we take
its 50 earliest carriers (sorted by filing year, conditional on top-topic
assignment with weight $\geq 0.20$), look up each owner's first granted mark
across the four classes, and classify the owner as \emph{entrant} if its first
grant is at the carrier-filing year (or one year earlier) versus
\emph{incumbent} otherwise.

\begin{table}[h]\centering\small
\begin{tabular}{lr}
\toprule
Themes scored                                    & 50 \\
Mean entrant share among 50 earliest carriers   & $73.7\%$ \\
Median                                            & $73.0\%$ \\
Themes entrant-dominated ($\geq 50\%$ entrants) & 48 / 50 \\
Themes incumbent-dominated ($\leq 20\%$ entrants) & 0 / 50 \\
\bottomrule
\end{tabular}
\end{table}

\noindent The raw numbers read as strong Schumpeter Mark~I --- 48/50 themes
entrant-dominated. \textbf{But the year-matched population baseline reverses
the claim.} Across the 4-class corpus the overall debut rate among granted
filings in 1990--2024 is \textbf{68--77\% per year} (mean $76.4\%$, weighted by
the year distribution of each theme's earliest carriers); the typical filer is
making their first grant. Theme-by-theme:

\begin{table}[h]\centering\small
\caption{RQ4 with population year-matched baseline.}
\begin{tabular}{lr}
\toprule
Theme entrant share (mean / median)              & 73.7\% / 73.0\% \\
Year-matched baseline (mean / median)            & 76.4\% / 76.7\% \\
\textbf{Excess (theme $-$ baseline)} & $\mathbf{-2.7}$pp / $-3.7$pp \\
Themes with positive excess                       & 22 / 50 \\
Themes with excess $>$ 10pp                       & 9 / 50 \\
Themes with negative excess                       & 28 / 50 \\
\bottomrule
\end{tabular}
\end{table}

\noindent The raw Schumpeter Mark~I read \textbf{is artifactual}: themes
inherit the corpus's high overall debut rate without adding to it. The
distribution of excesses is roughly centered on zero with a small entrant-lean
tail (9 themes at $+10$pp). \textbf{What this licenses}: \emph{the typical
theme is not disproportionately entrant-born once we account for who files in
the relevant years} --- a result worth reporting on its own. The 9 themes with
real excess merit a closer look (likely the most disruptive, e.g.,
blockchain/AI-era themes) and are the appropriate locus for any localized
Schumpeter~I claim. \textbf{Methodological note for the program:} this RQ
turned on the base rate, which the original §2.6 design did not specify ---
update the spec to require year-matched baseline for any future
provenance-typology result.

\section{Cross-method stability via NMF (§2.4 firewall)}
A 50-topic NMF was fit on the same 200{,}000-document stratified sample as the
LDA, using the same vocabulary ($V = 69{,}361$, $\min\,df = 200$), and
transformed across the full pool. Three comparisons:

\paragraph{Topic-word alignment (Jaccard on top-10 words, greedy match).}
Median Jaccard $= 0.18$, mean $= 0.24$; \textbf{28\% of NMF topics align with
an LDA topic at Jaccard $\geq 0.3$, 14\% at $\geq 0.5$}. The strongly-stable
intersection includes textbook business-model components:

\begin{table}[h]\centering\small
\caption{Top NMF$\leftrightarrow$LDA aligned topics.}
\begin{tabular}{ccp{0.62\textwidth}}
\toprule
J & pair & top words \\
\midrule
$1.00$ & NMF[9]/LDA[16] & energy, electricity, power, renewable \\
$0.82$ & NMF[14]/LDA[17]& bags, clothing, leather, cases, shoes \\
$0.67$ & NMF[0]/LDA[2]  & services, store, retail, wholesale \\
$0.67$ & NMF[5]/LDA[7]  & music, entertainment, video, audio \\
$0.67$ & NMF[17]/LDA[28]& systems, control, sensors, devices \\
\bottomrule
\end{tabular}
\end{table}

\noindent The bottom of the alignment is structural: NMF cleanly separates
\emph{networks/transmission} (NMF[44]) and \emph{cloud computing} (NMF[2])
where LDA blended cloud computing with lighting/fixtures (LDA[12]); NMF
isolates a tight VR/NFT/blockchain topic (NMF[45]) where LDA scattered it.
\textbf{The §2.4 firewall is therefore not idle: only the $\sim 25\%$ of LDA
topics with a Jaccard $\geq 0.3$ NMF counterpart should be promoted to
headline; the rest are method-artifacts to flag.}

\paragraph{Per-filing top-topic agreement.} Adjusted Rand Index $=+0.214$ on a
500{,}000-filing sample of paired (LDA, NMF) top-topic labels --- modest but
well above random ($\approx 0$). Methods partition the corpus similarly on
the stable themes and diverge on the rest, consistent with the Jaccard
distribution.

\paragraph{Communication with existing constructs.} For nine of ten diff2 seed
phrases there is a clear best-NMF-topic:

\begin{table}[h]\centering\small
\caption{Seed-phrase $\to$ NMF topic (best-match by phrase presence in top 50 words).}
\begin{tabular}{ll}
\toprule
seed phrase (diff2) & NMF topic (top-50 words contain) \\
\midrule
blockchain    & NMF[19]: cryptocurrency, blockchain \\
cloud         & NMF[2]: cloud computing \\
AI / ML       & NMF[16]: artificial intelligence, machine learning \\
SaaS / `as-a-service' & NMF[12]: software service \\
mobile-app    & NMF[48]: mobile applications \\
AR / VR       & NMF[45]: virtual reality, augmented reality, marketplace \\
marketplace   & NMF[45] (same as AR/VR --- a blended topic) \\
streaming     & --- no clean NMF topic \\
IoT           & --- no clean NMF topic \\
on-demand     & --- no clean NMF topic \\
\bottomrule
\end{tabular}
\end{table}

\noindent NMF \emph{independently rediscovers} seven of the ten hand-curated
seed phrases as topics. The three missing (streaming, IoT, on-demand) are
either too generic to be distinguished from other topics' top words or fall
below the cluster threshold at $T = 50$ --- worth re-checking at higher $T$.

\paragraph{Kill condition on NMF: pass, even more decisively than LDA.}
Re-running the shuffle null on the NMF top-topic panel: real mean lag-1
autocorr $= \mathbf{+0.689}$ (200 cells, essentially identical to LDA's
$+0.687$); shuffle baseline $= -0.021 \pm 0.012$ over three permutations;
\textbf{$z = +60.7$} (LDA was $+46$). The diffusion structure is
\emph{method-invariant} --- the trajectories are properties of the corpus, not
artifacts of either topic model.

\section{RQ5 --- margin vs.\ revenue demand-pull (associational)}
Using the SEC FSDS firm-year financial panel (2009--2024; 83{,}095 firm-year
rows, 13{,}127 unique CIKs) we link USPTO owners to CIKs via the exact-match
crosswalk (7{,}233 matched owners, 4{,}999 distinct CIKs in the borrowed-
novelty panel, 6{,}628 CIKs with a primary NICE class). For each (class,
year) we compute the class-level median gross margin and mean log-revenue
growth across firms whose primary NICE class falls in our subset. The
outcome --- theme inflow --- is operationalized two ways: $\Delta$ mean
share\_top and growth in the count of themes above the 1\% floor in the
class. We regress on lagged margin and lagged revenue growth with class and
year fixed effects.

\begin{table}[h]\centering\small
\caption{Class-level demand-pull regression ($n = 57$, $R^2 = 0.44$).}
\begin{tabular}{lrr}
\toprule
& coef & $t$ ($p$) \\
\midrule
lagged class margin             & $-11.34$ & $-1.41$ (0.16) \\
\textbf{lagged class revenue growth} & $\mathbf{+9.22}$ & $\mathbf{+1.98\ (0.048)}$ \\
\bottomrule
\end{tabular}
\end{table}

\noindent The mean-share-inflow outcome collapses numerically (themes are too
sparse for cell-mean changes to read above noise). On the count-of-active-
themes outcome, \textbf{lagged class revenue growth predicts theme uptake;
lagged margin does not.} Direction: \emph{Schmookler / demand-pull}
(invention follows market size), not appropriability/rents. \textbf{Caveat
carried openly:} with 4 NICE classes $\times$ $\sim$14 years $= 57$ degrees of
freedom and class+year FE, this is barely identified and labeled
associational. A 45-class build with NICE$\to$NAICS crosswalk would give the
proper test; the shock-based causal design (§4.7) is still owed.

\section{RQ6 --- does borrowed novelty pay? (firm-level outcome bridge)}
For the 4{,}999 CIK-linked firms in the borrowed-novelty panel (21{,}150
firm--filing-year rows), we aggregated per-filing borrowed-gap to firm-year
means and joined to SEC firm-year financials at lag 0, +1, and +2 years. We
standardize the firm-year borrowed gap to $\sigma$-units and fit firm + year
two-way FE with firm-clustered standard errors.

\begin{table}[h]\centering\small
\caption{Within-firm outcomes per $\sigma$ of borrowed gap.}
\begin{tabular}{lrrr}
\toprule
& gross margin & rev.\ growth & R\&D intensity \\
\midrule
contemp ($t$)   & $\mathbf{+0.014\ (t=+2.1)}$ & $+0.011\ (0.7)$ & $+0.086\ (0.4)$ \\
$t+1$           & $+0.007\ (0.8)$            & $-0.002\ (-0.1)$& $+0.085\ (0.3)$ \\
$t+2$           & $+0.013\ (1.6)$            & $-0.008\ (-0.6)$& $-0.142\ (-1.0)$ \\
\midrule
cross-section (firm means) & $+0.20\ (1.3)$  & $-0.003\ (-0.2)$ & $\mathbf{-15.5\ (-2.1)}$\textsuperscript{$\dagger$} \\
\bottomrule
\end{tabular}
\\[-2pt]
\textsuperscript{$\dagger$}\footnotesize unwinsorized, outlier-prone --- treat as suggestive.
\end{table}

\noindent Two findings worth keeping: \textbf{(a)} contemporaneous gross
margin moves up modestly with within-firm borrowed-novelty
($+1.4$pp per $\sigma$, $t = 2.1$), an effect that fades by $t+1$ ---
firms that import vocabulary into their class capture some same-year margin
lift but it does not persist as a multi-year premium. \textbf{(b)}
Cross-sectionally, high-borrowed-novelty firms have markedly \emph{lower}
R\&D intensity --- a \textbf{coherent picture with the construct-validity
finding (Part~I) that $\Delta$KL is a within-firm patent substitute}. Borrowed
novelty is the market-facing innovation channel; firms that pull this lever
spend less on the technical-invention channel. The two channels appear
substitutive within a firm, complementary across the population. This is
the localized firm-level innovation signal the original $\Delta$KL critique
demanded; the construct earns the bounded innovation reading on this
margin-and-substitution shape.

\section{What this licenses, and what is next}
\begin{itemize}[nosep]
\item The \textbf{descriptive diffusion claim is supported} (Result 1), the
\textbf{construct's gating tests pass} (Result 3), and a first-pass read of
RQ1--RQ4 (§2.5) produces internally consistent diffusion structure: ramped
adoption (autocorr $+0.69$), software-led cross-class flow, Bass-fittable
curves with plausible $q/p$, and a strikingly uniform entrant tilt
\emph{when read against the corpus base rate}.
\item The \textbf{theme is the unit}, not a per-filing scalar (Result 2 vs.\ 3
together).
\item \textbf{Phase 1 priorities, ranked by what the data now invites:}
RQ2 (breadth and class-to-class theme-flow network --- $T=50$ themes
$\times$ 45 classes once we widen the corpus); RQ3 (Bass $p/q/m$ fitted to the
transit themes' adoption curves --- the trajectories already look Bass-shaped
in Result 3); RQ1 (theme-birth Heaps' law); RQ4 (entrant vs.\ incumbent
provenance using the owner panel that is already on disk).
\item \textbf{Methodological work still owed before headline:} cross-method
theme stability (LDA \emph{and} embedding-cluster \emph{and} high-PMI
collocations, §2.4); human face-validation on a random theme sample (target
$\geq 0.8$ rater agreement); enriching the per-cell ahead-precedes-behind test
with length-weighting and within-theme-vintage controls (the per-filing signed
$\Delta\mathrm{KL}$ noise issue); a variable-$K$ theme model so RQ1's Heaps
reads beyond fixed-$T$ saturation; the §2.6 spec needs to fold in a baseline
requirement for any future provenance-typology claim.
\item \textbf{Now unblocked:} the SEC FSDS panel runs RQ5 (Schmookler-supportive
at the 4-class associational level) and RQ6 (modest contemporaneous margin
lift, R\&D substitute --- localized innovation signal on the bounded shape
the CV critique was willing to license). Both want a 45-class corpus and a
NICE$\to$NAICS crosswalk to read at full power.
\item \textbf{Open for human input:} \texttt{paper/face\_validation.md} ---
a 50-topic rateable form (top words, per-class peak, 8 representative marks,
checkboxes). Per §2.4, headline themes need rater agreement $\geq 0.8$ on the
``coherent'' bucket. Combined with the NMF stability filter (Jaccard $\geq 0.3$
counterpart, 14/50 strict / 28\% loose), the headline theme set is the
intersection of cross-method-stable AND human-validated. Currently
human-validation pending.
\end{itemize}

% ========================================================================
\chapter*{Reproducibility (Parts I and II)}
\addcontentsline{toc}{chapter}{Reproducibility (Parts I and II)}

\paragraph{Part~I (construct validity).} \texttt{scripts/wsC\_within\_firm\_patents.py},
\texttt{wsD\_validity\_battery.py}, \texttt{wsE\_returns\_diagnostic.py},
\texttt{wsF\_make\_figs.py}, plus class-009 corpus diagnostics
\texttt{wsG\_class009\_analysis.py}, \texttt{wsH\_term\_provenance009.py};
metrics in \texttt{paper/results/wsC\_*.json}, \texttt{wsD\_*.json},
\texttt{wsE\_*.json}, \texttt{wsG\_*.json}, \texttt{wsH\_*.json}.

\paragraph{Part~II (diffusion).} \texttt{scripts/diff1\_borrowed\_novelty.py},
\texttt{diff2\_theme\_transit.py}, \texttt{diff3\_lda\_themes.py},
\texttt{diff4\_phase1\_rqs.py}, \texttt{diff5\_rq4\_baseline.py},
\texttt{diff6\_nmf\_compare.py}, \texttt{diff7\_face\_validation\_doc.py},
\texttt{diff8\_rq6\_borrowed\_pays.py}, \texttt{diff9\_rq5\_demand\_pull.py};
JSON outputs in
\texttt{paper/results/\{diff1\_borrowed,diff2\_transit,diff3\_themes,diff4\_phase1,diff5\_rq4\_baseline,diff6\_nmf,diff8\_rq6,diff9\_rq5\}.json};
per-filing panels in
\texttt{data/processed/\{borrowed\_novelty,filing\_topics,filing\_topics\_nmf,theme\_panel,sec\_firm\_year,uspto\_sec\_crosswalk\}.parquet}.

\paragraph{Corpora.} Classes 009/042/039/035 from the USPTO TRTYRAP backfile
(91 zip parts, $\sim$12 GB), grant-conditioned, 1990--2024. SEC financial
panel: FSDS 2009Q1--2024Q4 (83{,}095 firm-years; 7{,}233 USPTO$\to$CIK exact
matches). Run all scripts under \texttt{.venv-analysis}
(numpy/pandas/scipy/statsmodels/scikit-learn/matplotlib).

\end{document}
